\documentclass[11pt]{article}
\usepackage{fontspec,amsmath,amssymb,amsthm,hyperref}
\usepackage[margin=1in]{geometry}
\setlength{\emergencystretch}{3em}
\title{The actual nonlinear kinetic coupling in L\ensuremath{{}^{2}}}
\author{LangevinLean mathematical exposition}
\date{}
\begin{document}\maketitle
\noindent Source module: \texttt{KineticLangevin/LpContraction.lean}. Compile this proof document with XeLaTeX or LuaLaTeX. The accompanying HTML pairs the mathematical argument with the exact Lean statements and records the scope of the formalization.


Random vectors form a Hilbert space under the L\ensuremath{{}^{2}} norm. The gradient acts on these vectors through the project's pointwise \texttt{lpDrift} construction. This module lifts the centered strong-smooth inequality to that actual vector field, then proves contraction of actual kinetic integral paths in L\ensuremath{{}^{2}}. It does not introduce a hypothetical differentiable potential on the infinite-dimensional space of random variables.


\section{Lift the centered gradient inequality pointwise}


For deterministic points, strong convexity and Lipschitz gradient imply \(\|\nabla f(x)-\nabla f(y)-\frac{m+M}{2}(x-y)\|\le\frac{M-m}{2}\|x-y\|\). The lifted gradient agrees almost everywhere with pointwise application of the original gradient. Lean collects the almost-everywhere identities for subtraction, scalar multiplication and \texttt{lpDrift\_coe}, rewrites the random-variable representatives, and applies the deterministic centered inequality at each sample. \texttt{Lp.norm\_le\_mul\_norm\_of\_ae\_le\_mul} integrates this norm comparison, yielding the same constant in L\ensuremath{{}^{2}}. No independence or moment estimate beyond membership in L\ensuremath{{}^{2}} is needed for this step.



\noindent Lean declarations: \texttt{lpDrift\_centered\_bound}.

\section{Derive the L\ensuremath{{}^{2}} energy inequality from actual integral paths}


Two phase paths satisfy the concrete kinetic integral equations with the actual lifted gradient and common additive input. The noise cancels in those equations, so their difference has a right derivative. Continuous linear projections extract the velocity and position equations. The general centered dissipativity theorem applies in the Hilbert space L\ensuremath{{}^{2}} because the preceding pointwise lift supplies its sector premise. Under \(\gamma^2\ge m+M\), the maximum curvature coefficient equals \(-m\), giving \(E'(t)\le-2mE(t)/\gamma\). Continuity and the right-slope Grönwall theorem prove \(E(t)\le E(0)e^{-2mt/\gamma}\). Taking nonnegative square roots, using the exponential-addition identity, gives the full transformed L\ensuremath{{}^{2}} norm contraction.



\noindent Lean declarations: \texttt{lp\_kinetic\_integral\_energy\_decay}, \texttt{lp\_kinetic\_integral\_transformNorm\_decay}.

\end{document}
